\Needspace{13\baselineskip}
\section{R6: Contracts, inductive proofs, and proof debt}\label{sec:contracts}

\Q{R6.Q1}{Interleaving programs and proofs}{What discipline avoids both premature proof attempts and discovering too late that a chosen value cannot satisfy its contract? Is there a Lean-native analogue of refinement abstraction?}

\answer Use a dependency-driven scheduler over both value and proof obligations, with cheap propagation before expensive proof search. The proposal's assertion that proving a property of a program with holes is not meaningful until those holes are filled is too strong. Reflexivity can prove \code{?n = ?n}; an equality proof obligation can constrain a witness; and a theorem quantified over the remaining holes can certify an entire family of completions. What is dangerous is treating a proof about one temporary assignment as valid for all later assignments.

There are three legitimate ways a proof obligation can make progress. It may be solved uniformly for all unresolved dependencies. It may assign some of those dependencies, committing those choices within the current branch. Or it may produce a conditional lemma whose remaining assumptions become explicit obligations. A proof scheduler should record which of these happened. A tactic that merely times out has supplied none of them.

For a pointwise postcondition $Q(x,y)$, search for
\[
 g : \prod_x \{y:B(x)\mid Q(x,y)\},
 \qquad f(x):=(g(x)).\mathsf{val}.
\]
An inductive branch receives both the recursive result and its property, making many obligations local. This is a useful form of correct-by-construction synthesis. It does not automatically decompose arbitrary relational contracts. An involution law $f(f(x))=x$, for example, couples different executions of $f$; a homomorphism law couples several inputs. Such contracts require a stronger invariant, a suitable algebraic construction, or a global proof.

Represent an obligation by its dependency set, a readiness class, previous proof attempts, and residual facts. Run inexpensive definitional equality, reflexivity, constructor reasoning, and already applicable lemmas early. Run a more expensive arithmetic or induction tier when relevant value structure stabilizes or when the obligation is sufficiently constraining to justify its cost. The policy must allow exceptions to a rigid ``all values first'' order: a proof about an index may be exactly what determines a value hole.

A finite predicate abstraction can be implemented without an external SMT solver. Choose a finite, explicit vocabulary of predicates on the current values, and use Lean-proved implication and transfer lemmas to propagate abstract facts. For example, summaries may record length equalities, nonemptiness, or interval bounds. A failed implication attempt means \emph{unknown}, not that the implication is false. Joins must forget unsupported facts rather than invent them. This is an engineering adaptation of refinement-style reasoning; it is not a complete abstraction of arbitrary Lean propositions. Synquid supplies a useful comparison point, but its solver-supported refinement fragment should not be conflated with unrestricted dependent proof search.\cite{synquid}

\Q{R6.Q2}{A recognizable class of automatic inductive contracts}{Which properties are reliably handled by induction along a synthesized program, simplification, and arithmetic? Can the engine identify that class?}

\answer Specify a sufficient, certificate-producing class rather than claiming to recognize every property that happens to admit such a proof. For a list fold with seed $z:K$, step $s:A\to K\to K$, and an invariant $I:\List A\to K\to\mathsf{Prop}$, require local certificates
\begin{align}
 &I([],z),\label{eq:inv-base}\\
 &\forall a\ xs\ k,\ I(xs,k)\to I(a::xs,s(a,k)).\label{eq:inv-step}
\end{align}
These imply the desired global fold invariant.

\begin{proposition}[Local algebra certificates]
If \cref{eq:inv-base,eq:inv-step} hold, then
\[
 \forall xs,\ I(xs,\fold\ s\ z\ xs).
\]
If additionally $I(xs,k)\to Q(xs,o(k))$, then the program $o\circ\fold\ s\ z$ satisfies $\forall xs, Q(xs,f(xs))$.
\end{proposition}
\begin{proof}
Induct on $xs$. The empty case is \cref{eq:inv-base}. In the constructor case, the induction hypothesis establishes $I(xs,\fold\ s\ z\ xs)$; applying \cref{eq:inv-step} gives the required invariant for $a::xs$. Apply the finishing implication to obtain the postcondition.
\end{proof}

This pattern extends to polynomial inductive datatypes by giving one preservation obligation per constructor and one invariant hypothesis per recursive field. For indexed families, the predicate and recursive results must be typed at their respective indices. The extension is structural, but the generated obligations can still be difficult.

A recognizable \emph{automation fragment} consists of annotated fold/recursor shapes, an invariant chosen from a fixed grammar, and local obligations which, after a terminating certified normalization, belong to a specified decidable theory. A useful initial theory is quantifier-free linear arithmetic over naturals and integers together with a fixed collection of constructor and projection facts. Universal branch variables become arbitrary local parameters; the resulting implication is discharged by the selected decision procedure. Completeness is then relative to the invariant grammar, normalization rules, and decision theory. It is not completeness for finding every inductive invariant.

For map-length preservation, choose $I(xs,ys)$ to be $\mathsf{length}(ys)=\mathsf{length}(xs)$. The empty case is immediate, and the cons case reduces to the recursive equality with one added to each side. An accumulator reversal needs a strengthened invariant involving the extra parameter; a property about sortedness or permutation needs corresponding logical structure and lemmas, not merely linear arithmetic. A syntactic recognizer should distinguish ``all local obligations lie in our certified fragment'' from ``the invariant is plausible but a general prover is needed.''

This also clarifies the role of helper invention. Inventing a pair carrier without an invariant can make the program search easier while making certification harder. Synthesize the carrier, its local invariant, and the finishing implication as related but separately budgeted objects. Do not charge a proof failure against the semantic correctness of the program unless a checked counterexample is obtained.

\Q{R6.Q3}{Available induction automation in Lean}{Which current Lean facilities can be used as bounded proof services, and what remains to be built?}

\answer The dependable core is an explicit induction/case-analysis adapter followed by ordinary proof-producing automation, not an assumption that a general-purpose tactic will invent the induction. Lean exposes the induction API discussed in R4; its tactic ecosystem provides simplification, arithmetic automation, and \code{grind}. Aesop's own documentation explicitly says that induction is deliberately left to the user in its standard approach. Canonical demonstrates relevant dependent inhabitation and recursor search, but that is not a promise that unrestricted program invariants will be discovered automatically.\cite{leaninduction,leanmanual,grind,aesop,canonical}

For Leant2, implement a small proof service which reads the synthesized program's recursion metadata, chooses its major argument, reverts the recorded generalized parameters, performs induction, unfolds the corresponding equation, and passes the resulting local goals to a tiered portfolio. The same metadata used to build the program should guide the proof; reconstructing it from a large opaque recursor expression is avoidable work. For a well-founded definition, use its checked equation and induction principle rather than pretending the body is definitionally structural.

LeanHammer is a relevant optional integration for premise selection and proof reconstruction; the premise-selection paper has an ICLR 2026 version, so the proposal's preliminary bibliography should be updated. It does not remove the need to identify the induction or to validate all reconstructed dependencies against the requested axiom policy. Its external infrastructure also makes it a different deployment profile from a Lean-core-only default.\cite{leanhammer}

Rippling and lemma discovery suggest useful search heuristics, but this investigation did not establish a maintained, drop-in Lean 4.34 library that supplies the entire proposed automatic-induction service. The practical answer is therefore a small native adapter plus explicitly optional integrations. A proof service should return a closed proof, a checked counterexample, a suspended obligation, or an explicit resource failure. It should never collapse these outcomes into a single Boolean.

Budgeting must cover the actual expensive operation. A deadline checked only between tactic calls cannot preempt a single long call. Use supported heartbeat/interrupt mechanisms inside the worker, and distinguish soft cooperative cancellation from a hard process boundary where one is required. Preserve the frozen environment and reconstruct the proof in the acceptance process. Millisecond tiers may be useful measured defaults on one machine; they are not reproducible semantic bounds.

\Q{R6.Q4}{Scheduling proof debt}{Is there a principled policy for spending time on proof obligations rather than more program search?}

\answer There is a principled engineering formulation, but no justified claim of a universally optimal bandit policy for this workload. Candidate branches share subterms, learned facts, and proof obligations; costs and success probabilities change as other goals are solved. This violates the simple independence and stationarity assumptions behind many elementary bandit models.

Start with a deterministic, fair policy. Cheap propagation runs at each relevant change. Expensive proof attempts receive increasing logical-work quanta, with separate queues for local branch certificates, termination proofs, and global postconditions. Reserve a nonzero share for older obligations so that promising new candidates cannot starve a correct existing one. Reuse a proof only when its dependency key still matches. A proof solved once for a generalized helper can be more valuable than several isolated global proof attempts.

An optional priority estimate is
\[
 \frac{\Pr(\text{certification or useful refutation}\mid\text{features})
       \cdot\text{expected downstream work saved}}
      {\text{estimated additional work}+\text{switching overhead}}.
\]
This is a scheduling heuristic, not a theorem about optimality. Learn it only after logging enough matched obligations to estimate both successes and unsuccessful costs. Freeze or record the estimator for replay; keep the certificate-producing services independent of it.

Checked counterexamples should be shared across candidates for the same frozen contract, including all necessary preconditions on their inputs. Survivors are tested on those counterexamples before expensive proofs. By contrast, ``the prover failed on this program'' must not become a shared negative example. The distinction preserves the rule that tests can refute a universal claim at a concrete point, while only a proof certifies the universal claim.
